\documentclass[11pt]{article}
\usepackage[margin=1in]{geometry}
\usepackage{amsmath,amssymb,amsfonts}
\usepackage{braket}
\usepackage{microtype}
\usepackage[hidelinks]{hyperref}
\setlength{\parindent}{0pt}
\setlength{\parskip}{0.55em}
\title{\textbf{PHYS 3316 -- Quantum Physics I}\\
Class Notes -- 22 September 2026}
\author{Eric R. Bittner}
\date{}
\begin{document}
\maketitle
\section{Particle Subject to a Constant Force}
Consider a particle subject to a constant force, with a linear potential.
The Hamiltonian is
\[
H=\frac{p^2}{2m}+V(x).
\]
From Ehrenfest's theorem,
\[
i\hbar\frac{d}{dt}\langle x\rangle
=
\langle[x,H]\rangle
=
\frac{i\hbar}{m}\langle p\rangle,
\]
so that
\[
\boxed{
\frac{d\langle x\rangle}{dt}
=
\frac{\langle p\rangle}{m}.
}
\]
Likewise,
\[
i\hbar\frac{d}{dt}\langle p\rangle
=
\langle[p,H]\rangle
=
\frac{\hbar}{i}
\left\langle\frac{\partial V}{\partial x}\right\rangle.
\]
Thus
\[
\boxed{
\frac{d\langle p\rangle}{dt}
=
-\left\langle\frac{\partial V}{\partial x}\right\rangle.
}
\]
For a constant force $f$ we may write
\[
\frac{d\langle p\rangle}{dt}=f,
\]
and hence
\[
\langle p(t)\rangle
=
\langle p(0)\rangle+ft.
\]
If
\[
\langle p(0)\rangle=0,
\]
then
\[
\langle x(t)\rangle
=
\frac{f}{2m}t^2.
\]
% The handwritten notes use both V(x)=fx and a force denoted +f.
% Since F=-dV/dx, this corresponds to a sign-convention change between
% portions of the board work. The equations above preserve the stated
% constant-force equation d<p>/dt=f.
\subsection{Momentum Dispersion}
Consider
\[
\frac{d}{dt}
\left\langle
(p-\langle p\rangle)^2
\right\rangle
=
\frac{d\langle p^2\rangle}{dt}
-
\frac{d\langle p\rangle^2}{dt}.
\]
Using Ehrenfest's theorem,
\[
i\hbar\frac{d}{dt}\langle p^2\rangle
=
\langle[p^2,H]\rangle.
\]
For a linear potential,
\[
[p^2,V]
=
p[p,V]+[p,V]p,
\]
so that
\[
\frac{d}{dt}\langle p^2\rangle
=
2f\langle p\rangle.
\]
Also,
\[
\frac{d}{dt}\langle p\rangle^2
=
2\langle p\rangle
\frac{d\langle p\rangle}{dt}
=
2f\langle p\rangle.
\]
Therefore
\[
\boxed{
\frac{d}{dt}(\Delta p)^2=0.
}
\]
Thus the momentum distribution translates under a constant force without
changing its width: the motion is dispersionless in momentum space.
\section{Momentum-Space Solution}
Now solve the same constant-force problem directly in momentum space. In
position space,
\[
i\hbar\frac{\partial}{\partial t}\psi(x,t)
=
\left[
-\frac{\hbar^2}{2m}\frac{\partial^2}{\partial x^2}
+
fx
\right]\psi(x,t).
\]
Under the Fourier transform,
\[
x\longrightarrow i\hbar\frac{\partial}{\partial p},
\qquad
-i\hbar\frac{\partial}{\partial x}\longrightarrow p.
\]
Therefore
\[
i\hbar\frac{\partial\widetilde\psi}{\partial t}
=
\left[
\frac{p^2}{2m}
+
i\hbar f\frac{\partial}{\partial p}
\right]\widetilde\psi,
\]
or
\[
\frac{\partial\widetilde\psi}{\partial t}
=
f\frac{\partial\widetilde\psi}{\partial p}
-
\frac{i}{\hbar}\frac{p^2}{2m}\widetilde\psi.
\]
This is a first-order transport equation in $p$. Its characteristics obey
\[
\frac{dp}{dt}=-f,
\]
which is the classical force law for the sign convention $V(x)=fx$:
\[
-\frac{\partial V}{\partial x}=-f.
\]
Along a characteristic,
\[
p(t)=p_0-ft,
\]
and
\[
\frac{d\widetilde\psi}{dt}
=
-\frac{i}{2m\hbar}p(t)^2\widetilde\psi.
\]
Thus, if
\[
\widetilde\psi(p,0)=\widetilde\psi_0(p),
\]
the solution may be written
\[
\boxed{
\widetilde\psi(p,t)
=
\widetilde\psi_0(p+ft)
\exp\left[
-\frac{i}{2m\hbar}
\left(
p^2t+pft^2+\frac{f^2t^3}{3}
\right)
\right].
}
\]
The momentum distribution therefore simply translates while accumulating a
phase.
\section{Time Dependence of a Free Gaussian Wavepacket}
Consider a free particle initially localized in a Gaussian state,
\[
\psi(x,0)
=
A e^{-\sigma x^2},
\]
with normalization
\[
A=
\left(\frac{2\sigma}{\pi}\right)^{1/4}.
\]
Transform to momentum space:
\[
\phi(p)
=
\frac{1}{\sqrt{2\pi\hbar}}
\int_{-\infty}^{\infty}
e^{-ipx/\hbar}\psi(x)\,dx.
\]
For the Gaussian,
\[
\phi(p)
=
\frac{A}{\sqrt{2\sigma\hbar}}
\exp\left(
-\frac{p^2}{4\sigma\hbar^2}
\right).
\]
For a free particle,
\[
H=\frac{p^2}{2m},
\]
so
\[
i\hbar\frac{\partial\phi}{\partial t}
=
\frac{p^2}{2m}\phi.
\]
The momentum-space wavefunction therefore evolves only by a phase:
\[
\boxed{
\phi(p,t)
=
\phi(p,0)
\exp\left(
-\frac{ip^2t}{2m\hbar}
\right).
}
\]
In particular,
\[
|\phi(p,t)|^2=|\phi(p,0)|^2;
\]
the momentum distribution does not change with time.
Transforming back to position space,
\[
\psi(x,t)
=
\frac{1}{\sqrt{2\pi\hbar}}
\int_{-\infty}^{\infty}
\phi(p,t)e^{ipx/\hbar}\,dp.
\]
The remaining integral is again Gaussian and gives
\[
\boxed{
\psi(x,t)
=
A\,
\frac{1}{\sqrt{1+2i\hbar\sigma t/m}}
\exp\left[
-\frac{\sigma x^2}
{1+2i\hbar\sigma t/m}
\right].
}
\]
Define the characteristic spreading time
\[
\boxed{
\tau=\frac{m}{2\hbar\sigma}.
}
\]
Then
\[
\psi(x,t)
=
A\,
\frac{1}{\sqrt{1+it/\tau}}
\exp\left[
-\frac{\sigma x^2}{1+it/\tau}
\right].
\]
The probability density remains Gaussian, but its spatial width increases
with time. Initially,
\[
(\Delta x)_0^2=\frac{1}{4\sigma},
\]
while the momentum variance remains fixed:
\[
(\Delta p)^2=\sigma\hbar^2.
\]
The spatial width is
\[
\boxed{
\Delta x(t)
=
\frac{1}{2\sqrt{\sigma}}
\sqrt{
1+
\left(
\frac{2\hbar\sigma t}{m}
\right)^2
}.
}
\]
Equivalently,
\[
\boxed{
(\Delta x(t))^2
=
(\Delta x_0)^2
+
\left(
\frac{\Delta p}{m}t
\right)^2.
}
\]
The second term represents the spreading caused by the distribution of
velocities present in a localized state.
If the initial packet has a nonzero mean momentum,
\[
\psi(x,0)
=
A e^{-\sigma x^2}e^{ip_0x/\hbar},
\]
then
\[
\langle v\rangle=\frac{p_0}{m},
\]
while the packet undergoes the same spreading described above.
\end{document}
